\subsection{Реализация MFG и NMFG}
\subsubsection{решение (псевдокод)}

Алгоритм вычисления равновесия в MFG решается итеративно с помощью пары уравнений Гамильтона–Якоби–Беллмана (ГЯБ) и Фоккера–Планка (ФП). В сущности, мы начинаем с начального распределения плотности $\rho^0(x)$ и повторно выполняем следующие шаги до \cite{cardaliaguet2013notes}:

\begin{algorithm}[H]
\caption{Итеративное вычисление равновесия для MFG/NMFG}
\begin{algorithmic}[1]
    \State \textbf{Шаг 1:} Решить уравнение ГЯБ (обратное по времени) для функции стоимости $V(x,t)$ при данном распределении $\rho$.
    \State \textbf{Шаг 2:} На основе полученного $V$ вычислить оптимальную политику (управление) $\alpha(x,t)$, максимизирующую выигрыш при текущей $\rho$.
    \State \textbf{Шаг 3:} Решить уравнение ФП (прямое по времени) для $\rho^{\text{new}}$, используя управление $\alpha$.
    \State \textbf{Шаг 4:} Если норма разности $\|\rho^{\text{new}} - \rho\|$ стала меньше порога, завершить итерации; иначе присвоить $\rho \gets \rho^{\text{new}}$ и повторить.
\end{algorithmic}
\end{algorithm}

Реализация для NMFG (многоклассовой системы) аналогична, но с отдельным распределением $\rho_i$ для каждого класса $i$. То есть на каждом шаге одновременно решаются несколько пар уравнений (ГЯБ и ФП) — по одной паре на каждый класс — с учетом взаимодействия распределений разных классов. 

Временная сложность такого алгоритма оценивается как $O(K \cdot N^2 \cdot T)$ на итерацию (решение двухмерной задачи ГЯБ–ФП на сетке размером $N \times N$ при $T$ временных слоях и $K$ итерациях). Для двух классов NMFG это примерно в 2 раза медленнее, так как требуется вдвое больше уравнений, что соответствует оценке $O(2K \cdot N^2 \cdot T)$.

\subsubsection{решение (код)}

код на github будет скорее всего (чтобы меньше места занимать)

\subsubsection{мини заключение: сравнение ожидаемого времени и полученного экпеременально (или мб ещё что-то)}

Экспериментально подтверждается, что алгоритм NMFG работает примерно вдвое медленнее одноклассового MFG. На сетке $100 \times 100$ при той же точности сходимости время решения NMFG получалось примерно в $1.5\text{--}2$ раза больше, что соответствует оценке $O(2K \cdot N^2)$ против $O(K \cdot N^2)$ для MFG. При этом реальное время зависит от деталей реализации (например, выбора схемы конечных разностей) и типа оборудования. Замечено, что итерационный метод (решение ГЯБ и ФП) действительно является вычислительно дорогим процессом, но разделение на классы (NMFG) приводит к линейному росту затрат по числу классов.
